\section{Maximally entangled states and Choi matrices}
\label{sec:choi_jamiolkowski}

The maximally entangled state and the flip operator give the basic matrices
used by the Choi--Jamio\l{}kowski correspondence.

\begin{definition}[Maximally entangled state]\label{def:maximally_entangled}
    \lean{Matrix.omegaProj}
    \leanok
    For every $d\in\N$, define $\ketbra{\Omega}{\Omega}$ on
    $\C^d\otimes\C^d$ using
    \begin{align}
        \ket{\Omega}
        &= \frac{1}{\sqrt{d}} \sum_{j=0}^{d-1} \ket{j,j}.
        \label{eq:representations_max_entangled}
    \end{align}
    When $d\geq1$, this is the \emph{maximally entangled state}: a rank-one
    projector with
    $(\ketbra{\Omega}{\Omega})_{(i_1,i_2),(j_1,j_2)}
        = \frac{1}{d}\delta_{i_1 i_2}\delta_{j_1 j_2}$
    and $\tr(\ketbra{\Omega}{\Omega})=1$.
\end{definition}

\begin{definition}[Flip operator]\label{def:flip_operator}
    \lean{Matrix.swapMatrix}
    \leanok
    The \emph{flip operator} $F$ on $\C^d \otimes \C^d$ is
    \begin{align}
        F
        &= \sum_{i,j=0}^{d-1} \ket{ij}\!\bra{ji}.
        \label{eq:representations_flip}
    \end{align}
    Its matrix entries are
    \begin{align}
        F_{(i_1,i_2),(j_1,j_2)}
        &=
        \begin{cases}
            1, & i_1=j_2 \text{ and } i_2=j_1,\\
            0, & \text{otherwise}.
        \end{cases}
        \notag
    \end{align}
\end{definition}

\begin{definition}[Two-dimensional antisymmetric vector]
    \label{def:fin_two_antisymm_vec}
    \lean{Matrix.finTwoAntisymmVec}
    \leanok
    In $\C^2 \otimes \C^2$, let $v = \ket{01} - \ket{10}$.
    Equivalently, $v_{(0,1)}=1$, $v_{(1,0)}=-1$, and all other coefficients
    are zero.
\end{definition}

\begin{definition}[Tensor product of a map with identity]\label{def:tensor_map_id}
    \lean{Matrix.tensorMapId}
    \leanok
    Given a linear map $T : \MN{D} \to \MN{D}$,
    the \emph{tensor extension} $T \otimes \id$ acts on bipartite matrices
    $X \in \MN{D \times D}$ by applying $T$ to each ``slice'':
    \begin{align}
        (T \otimes \id)(X)_{(i_1,i_2),(j_1,j_2)}
        &= (T(X^{(i_2,j_2)}))_{i_1 j_1},
        \label{eq:representations_tensor_map_id}
    \end{align}
    where $X^{(i_2,j_2)}_{ab} = X_{(a,i_2),(b,j_2)}$ is the
    bipartite slice.
\end{definition}

%% =========================================================

The Choi--Jamio\l{}kowski isomorphism bends the input legs of the channel $T$
around the maximally entangled pair $\ket{\Omega}$, presenting the map as the
matrix $\tau=(T\otimes\id)\ketbra{\Omega}{\Omega}$ on the doubled space:
\begin{center}
    % Formula: \tau=(T\otimes\id)(\ketbra{\Omega}{\Omega}).
    % Source verdict: Wolf, Proposition 2.1 and Equation (2.1), applies the
    %   two-wire map T\otimes\id to the maximally entangled projector.
    % In indices, \tau_{(i_1,i_2),(j_1,j_2)}
    %   =D^{-1}(T(E_{i_2j_2}))_{i_1j_1}; there is no summed Kraus index and
    %   therefore no pairleg between the two rails.
    % The east matrix cup is the input projector \ketbra{\Omega}{\Omega}.
    % The opaque two-wire atom acts by T on one tensor factor and by \id on
    %   the other; the two west rails are the output indices of \tau.
    % Boundary signature: (virtual-west=2 | virtual-east=0 |
    % physical-up=0 | physical-down=0).
    \begin{tenkz}[
        rows={wire,wire},
        east={cup=$\ketbra{\Omega}{\Omega}$},
        west label=$\tau$
    ]
        \tnghost{} & \tnX[wires=2]{T\otimes\id} & \tnghost{} \\
        \tnghost{} & & \tnghost{}
    \end{tenkz}
\end{center}

\begin{definition}[Choi matrix]\label{def:choi_matrix}
    \lean{ChoiJamiolkowski.choiMatrix}
    \leanok
    \uses{def:maximally_entangled, def:tensor_map_id}
    The \emph{Choi matrix} of a linear map
    $T : \MN{D} \to \MN{D}$ is
    \begin{align}
        \tau
        &= (T \otimes \id)(\ketbra{\Omega}{\Omega})
        \in \MN{D \times D}.
        \label{eq:representations_choi}
    \end{align}
    Concretely,
    $\tau_{(i_1,i_2),(j_1,j_2)}
        = \frac{1}{D}\,(T(E_{i_2 j_2}))_{i_1 j_1}$
    where $E_{i_2 j_2}$ is the matrix unit.
    This is the Choi--Jamiol\-kowski convention of
    \cite[Proposition~2.1]{Wolf2012Quantum}.

    % Scope restriction (d=d'): this definition treats only endomorphisms of
    % one matrix algebra.  The different-dimension correspondence is recorded
    % in docs/paper-gaps/choi_rectangular_scope.tex.
\end{definition}

\begin{theorem}[Choi matrix of the identity map]\label{thm:choi_matrix_id}
    \lean{ChoiJamiolkowski.choiMatrix_id}
    \leanok
    \uses{def:choi_matrix, def:maximally_entangled}
    The identity channel has Choi matrix
    $\ketbra{\Omega}{\Omega}$.
\end{theorem}

\begin{proof}\leanok
    Apply the identity map in the definition of the Choi matrix.
\end{proof}

\begin{theorem}[Choi matrix of transposition]\label{thm:choi_transpose_flip}
    \lean{ChoiJamiolkowski.choiMatrix_transposeLinearMapComplex_eq_swap}
    \leanok
    \uses{def:choi_matrix, def:flip_operator, def:maximally_entangled}
    Let $\theta(X)=X^T$ be matrix transposition on $\MN{D}$, with $D \ge 1$.
    Then
    \begin{align}
        (\theta\otimes\id)(\ketbra{\Omega}{\Omega})
        &= \frac{1}{D}F.
        \label{eq:representations_transpose_choi}
    \end{align}
    This is \cite[Equation~3.1]{Wolf2012Quantum}.
\end{theorem}

\begin{proof}\leanok
    \uses{def:choi_matrix, def:flip_operator}
    The $(i_2,j_2)$ slice of $\ketbra{\Omega}{\Omega}$ is
    $D^{-1}E_{i_2j_2}$. Transposition sends this matrix unit to
    $D^{-1}E_{j_2i_2}$, whose $(i_1,j_1)$ entry is $D^{-1}$ exactly when
    $i_1=j_2$ and $i_2=j_1$.
\end{proof}

\begin{theorem}[Complete positivity iff the Choi matrix is positive semidefinite]\label{thm:cp_iff_choi_posSemidef}
    \lean{ChoiJamiolkowski.cp_iff_choi_posSemidef}
    \leanok
    \uses{def:choi_matrix, def:cp_map}
    Let $D\geq1$. A linear map $T : \MN{D} \to \MN{D}$ is completely
    positive (in the Kraus sense) if and only if its Choi matrix $\tau \ge 0$.
    This is the $d=d'$ specialization of
    \cite[Proposition~2.1]{Wolf2012Quantum}.
    % Scope restriction (d=d'): the different-dimension correspondence is
    % recorded in docs/paper-gaps/choi_rectangular_scope.tex.
\end{theorem}

\begin{proof}
    \leanok
    \uses{def:choi_matrix, def:cp_map}
    ($\Rightarrow$) If $T(X) = \sum_i K_i X K_i^\dagger$, then
    $\tau = \sum_i \ketbra{v_i}{v_i}$ where
    $(v_i)_{(a,b)} = c\,K_i(a,b)$ and $c = 1/\sqrt{D}$.
    This is a sum of rank-one PSD matrices.

    ($\Leftarrow$) If $\tau \ge 0$, write
    $\tau = \sum_m \ketbra{v_m}{v_m}$
    by spectral decomposition, define $K_m(a,b) = v_m(a,b)/c$,
    and recover $T(X) = \sum_m K_m X K_m^\dagger$ by linearity on
    matrix units.
\end{proof}

Unitary changes of Kraus operators preserve the corresponding completely
positive map.

\begin{theorem}[Unitary freedom in Kraus operators]\label{thm:kraus_same_map_of_unitary_combination}
    \lean{kraus_same_map_of_unitary_combination}
    \leanok
    \uses{def:cp_map}
    Let $U$ be a unitary $r \times r$ matrix ($U^\dagger U = \Id$)
    and suppose $K_j = \sum_\ell U_{j\ell} \tilde{K}_\ell$.
    Then $\{K_j\}$ and $\{\tilde{K}_\ell\}$ define the same Kraus map:
    for every $X \in \MN{D}$,
    \begin{align}
        \sum_j K_j X K_j^\dagger
        &= \sum_\ell \tilde{K}_\ell X \tilde{K}_\ell^\dagger.
        \notag
    \end{align}
\end{theorem}

\begin{proof}
    \leanok
    Substitute the combination formula, expand the triple sum over
    $j,\ell,\ell'$, swap summation order, and use the entry-wise form
    of $U^\dagger U = \Id$ to collapse to $\ell = \ell'$ diagonal terms.
\end{proof}

\begin{theorem}[Rectangular Kraus freedom]\label{thm:kraus_rectangular_freedom}
    \lean{kraus_rectangular_freedom}
    \leanok
    \uses{def:cp_map}
    If $\sum_\alpha B_\alpha X B_\alpha^\dagger = \sum_j A_j X A_j^\dagger$
    for all $X$ and $r_2 \le r_1$ (where $\{B_\alpha\}_{\alpha=0}^{r_1-1}$
    is the larger family and $\{A_j\}_{j=0}^{r_2-1}$ the smaller),
    then there exists a rectangular isometry
    $V$ ($r_1 \times r_2$, $V^\dagger V = \Id_{r_2}$) such that
    $B_\alpha = \sum_j V_{\alpha j}\,A_j$.
\end{theorem}

\begin{proof}
    \leanok
    Pad the smaller family by zero operators to a family $A'$ of size $r_1$.
    Evaluating equality of the two maps on matrix units and taking matrix
    entries gives equality of the Gram matrices of the Stinespring vectors,
    $M_B^\dagger M_B=M_{A'}^\dagger M_{A'}$.  Hence the assignment from
    each column of $M_{A'}$ to the corresponding column of $M_B$ is
    well-defined and isometric on their common span.  Extend this isometry to
    the ambient space $\C^{r_1}$ and let $W$ be its unitary matrix.  Restricting
    $W$ to the original $r_2$ coordinates gives a rectangular matrix $V$ with
    $V^\dagger V=\Id_{r_2}$.  Reading the column identity entrywise yields
    $B_\alpha=\sum_jV_{\alpha j}A_j$.
\end{proof}
